\chapter{Introduction}
\label{ch:intro}

\section{Motivation}

A natural language autoencoder (NLA) takes one activation vector from inside a language model and writes a
short paragraph describing what that vector contains~\cite{frasertaliente2026nla}. A second model reads the
paragraph and tries to rebuild the vector. If the rebuilt vector is close to the original, the paragraph must
have carried most of the information, so it can be read as an explanation of the activation. The appeal is
obvious: instead of a list of feature numbers, you get English.

The trouble is that English is persuasive. Here is a passage from Wikipedia, and what the released NLA for
Qwen2.5-7B wrote about the activation one token after the word \emph{cello}:

\begin{sample}{Passage (the model has read up to ``son'')}
\ldots It has since performed at the Wigmore Hall in London, and in many countries worldwide. A recording of
the \true{cello} son$|$atas of Antonio Vivaldi \ldots
\end{sample}

\begin{sample}{Explanation written by the NLA}
Biographical music description format with British classical composer context \ldots\ Final token ``son'' is
a truncated word mid-expression (``the \wrong{piano} son''), part of a noun phrase listing the completed
work's content (``the \wrong{piano} sonatas''), strongly expecting ``atas'' \ldots
\end{sample}

Most of this is right. The explanation knows the text is about classical music in a British setting, it
knows the current word is the first half of ``sonatas'', and it knows which letters come next.
The one thing it gets wrong is the specific: the instrument was a cello, and the explanation says piano, twice,
with no hint of doubt. A reader who trusts the explanation walks away believing the model was thinking about a
piano.

This report is about that failure. I call the cello a \emph{specific}: a named entity that is really in the
text. The gist (classical music, a recording, sonatas) survives the round trip through language. The
specific often does not, and what replaces it is a confident wrong name of the same kind.

\section{Prior Work}

Wrong specifics in NLA explanations are a known problem, and several groups have worked on it.

\begin{itemize}
\item \emph{The NLA authors} report it themselves~\cite{frasertaliente2026nla}: explanations are usually right
about the theme and can be wrong about the detail, and a claim that recurs across neighbouring tokens is more
likely to be true.
\item \emph{Audits.} Dingeto~\cite{dingeto2026recap} finds that on the released Qwen2.5-7B verbalizer only
about 2\% of specific claims are ones the reconstruction depends on. Denamgana\"i~\cite{denamganai2026monitorability}
finds that a true account beats a rival differing in one claim only 60\% of the time.
\item \emph{Fixes by retraining.} Flow-NLA~\cite{lek2026flownla} replaces point reconstruction with a model of
all the activations an explanation is compatible with. Hong et al.~\cite{hong2026grounding} add a grounding
reward for names, numbers and quotations. Dingeto's RECAP~\cite{dingeto2026recap} trains the target model
itself so that chosen content stays readable.
\item \emph{Neighbouring tools.} Other ways of turning activations into text include the logit
lens~\cite{nostalgebraist2020logitlens}, Patchscopes~\cite{ghandeharioun2024patchscopes} and
SelfIE~\cite{chen2024selfie}; sparse autoencoders are the main
alternative~\cite{bricken2023monosemanticity,cunningham2023sparse}. Linear probes~\cite{alain2016probes} test
what an activation holds, and agreement between samples is a standard signal for invented content in
ordinary generation~\cite{manakul2023selfcheckgpt,wang2023selfconsistency}.
\end{itemize}

\section{Question}

I ask one narrow question and try to answer it carefully:

\begin{quote}
When an NLA explanation names an entity, how often is the name wrong, why does that happen, and what removes
the wrong names?
\end{quote}

The middle part matters most. The aim is to understand why the verbalizer behaves this way, not to ship a
better one. So \emph{how often} is a measurement against a baseline, \emph{why} is tested with an
intervention, and the fixes are there mainly as checks on the explanation: if wrong names are free under the
training objective, more training on that objective should not remove them, and a penalty on them should.

\section{Contributions}

Working with the released Qwen2.5-7B NLA and 6{,}000 Wikipedia passages about 60 entities, this report adds
four things to the work above. As far as I know each is new. None of them retrains an NLA from scratch or
needs the source text at inference.

\begin{enumerate}
\item \textbf{A fade curve with a ceiling.} Names are tracked token by token after the entity, next to a probe
on the same activations, which separates what the verbalizer drops from what the model no longer holds.
Sixteen tokens after the entity the explanation names it 7.5\% of the time while the probe still reads it
35\% of the time, and between a fifth and a half of all explanations contain a wrong name of the same type
(Chapter~\ref{ch:fading}).

\item \textbf{A price for each kind of name.} Sharpening Dingeto's 2\% finding, an edit test prices wrong and
true names separately. A wrong name is worth about 0.002 of reconstruction score. The true name is worth
0.21 at the entity's token and 0.016 one token later (Chapter~\ref{ch:incentives}).

\item \textbf{Agreement as a filter, and where it comes from.} The NLA authors' observation about neighbouring
tokens is applied to repeated samples of one activation and measured. A name in one of four samples is right
6\% of the time; a name in all four, 64\%. Keeping names that two samples agree on removes 74\% of the wrong
names and keeps 77\% of the true ones, with no training and no labels. Most of this signal is token-level
confidence that a single explanation already carries (Chapter~\ref{ch:agreement}).

\item \textbf{Rewards as tests.} Small fine-tunes of the released verbalizer check the incentive account. More
reconstruction leaves wrong names untouched. A cost on names the other samples do not support cuts them by
5 to 8\% on held-out passages, in two independent runs, without losing true names. Rewarding the true name
does not bring it back (Chapter~\ref{ch:finetuning}).
\end{enumerate}

Along the way I found that the usage example published with the models reads activations one layer too early,
which silently lowers reconstruction quality; Appendix~\ref{app:layer} documents the bug and the check that confirmed the
fix. Every table in this report is
generated from saved outputs by one command.

\paragraph{Availability.} Everything needed to check or rerun the work is in three places:
{\small
\begin{itemize}
\item Code, explanations and result tables:\\ \url{https://github.com/tanvincible/nla-wrong-names}
\item Fine-tuned adapters:\\ \url{https://huggingface.co/tanvincible/nla-qwen2.5-7b-L20-av-specifics-lora}
\item Activations, caches and training samples:\\ \url{https://huggingface.co/datasets/tanvincible/nla-specifics-artifacts}
\end{itemize}}
